% Student paper — English, typeset text, graphs as CIE clips. Compiler: XeLaTeX.
\documentclass[10pt,a4paper]{article}

\usepackage[margin=16mm,headheight=13pt,headsep=5pt,footskip=9mm]{geometry}
\usepackage{fontspec}
\usepackage[version=4]{mhchem}
\usepackage{chemfig}
\usepackage{graphicx}
\usepackage{xcolor}
\usepackage{booktabs}
\usepackage{array}
\usepackage{enumitem}
\usepackage{needspace}
\usepackage{fancyhdr}
\usepackage{amsmath,amssymb}
\usepackage{pgffor}

\raggedbottom
\setlist[enumerate]{leftmargin=1.5em,itemsep=0.12em,topsep=0.15em,parsep=0pt}
\setlist[itemize]{leftmargin=1.3em,itemsep=0.08em}

\pagestyle{fancy}
\fancyhf{}
\fancyhead[L]{\small \homeworktitle}
\fancyhead[R]{\small \homeworkmarks}
\fancyfoot[C]{\thepage}
\renewcommand{\headrulewidth}{0.3pt}

\newcommand{\homeworktitle}{Chemistry homework}
\newcommand{\homeworkmarks}{}
\newcommand{\sethomework}[2]{%
  \renewcommand{\homeworktitle}{#1}%
  \renewcommand{\homeworkmarks}{#2}%
}

\newcommand{\answerlines}[1]{%
  \par\vspace{0.08em}%
  \foreach \n in {1,...,#1}{\noindent\rule{\linewidth}{0.2pt}\\[0.62em]}%
}

\graphicspath{{images/}}
\DeclareGraphicsExtensions{.png,.jpg,.jpeg,.pdf}

\setlength{\parindent}{0pt}
\setlength{\parskip}{0.10em}
\setlength{\abovedisplayskip}{0.25em}
\setlength{\belowdisplayskip}{0.25em}

\newcommand{\qhead}[2]{%
  \par\vspace{0.28em}%
  \Needspace{4.2em}%
  \noindent{\small\textbf{#1}}\hfill{\small[#2]}\par\vspace{0.08em}%
}

% Graph clip: stem is on the scan. Measure so label + image stay together.
\newcommand{\clipq}[3]{%
  \par
  \sbox0{\includegraphics[width=0.84\linewidth,keepaspectratio]{#3}}%
  \Needspace{\dimexpr\ht0+1.25em+2pt\relax}%
  \noindent{\small\textbf{#1}}\hfill{\small[#2]}\par\vspace{0.06em}%
  \noindent\usebox0\par
}

\newcommand{\optline}[1]{%
  \par\vspace{0.06em}{\small #1}\par
}

\newcommand{\sourcefoot}[1]{%
  {\footnotesize\color{gray}#1}\par
}

\begin{document}
\sethomework{AS Chemistry homework · 3.1 Electronegativity}{16 marks}

\begin{center}
{\large\bfseries AS Chemistry homework · 3.1 Electronegativity}\\[0.12em]
{\small CIE 9701 AS Chemistry · 16 marks · about 35--40 minutes}
\end{center}

{\small Topic: 3.1 Electronegativity and bonding --- definition, factors, Periodic Table trends, and using Pauling differences to predict ionic or covalent bonding.
Answer in English. Section A is Paper 1 style (1 mark each). Section B is written; give the usual CIE reasons (nuclear charge, atomic radius, shielding). Pauling values are given on the paper where needed.}

\vspace{0.15em}
{\small
\begin{tabular}{@{}llr@{}}
\toprule
\textbf{Section} & \textbf{Content} & \textbf{Marks} \\
\midrule
A & Multiple choice & 11 \\
B1 & 2023 FM Paper 22 Question 1(a) & 3 \\
B2 & 2023 MJ Paper 22 Question 1(b)(ii) & 2 \\
\midrule
 & \textbf{Total} & \textbf{16} \\
\bottomrule
\end{tabular}}

\vspace{0.25em}
{\large\bfseries Section A --- Multiple-choice questions (11 marks)}\par

% A1 graphs --- keep CIE clip
\clipq{A1.}{1}{img-001.png}
\sourcefoot{2022 MJ Paper 12 Question 19}

% A2 NEW --- typeset
\qhead{A2.}{1}
{\small V and Z are both elements in Period 3 of the Periodic Table. Each element forms one stable ion that does not contain another element.}

{\small The atomic radius of each element and the ionic radius of the ion described above is shown.}

\vspace{0.15em}
\centering
{\small
\begin{tabular}{ccc}
\toprule
element & atomic radius / nm & ionic radius / nm \\
\midrule
V & 0.186 & 0.095 \\
Z & 0.099 & 0.181 \\
\bottomrule
\end{tabular}}
\par\raggedright

\vspace{0.12em}
{\small Which statement is correct?}
\begin{enumerate}[label=\textbf{\Alph*}, leftmargin=1.7em, itemsep=0.08em]
\item Ions of V and Z have the same number of full electron shells.
\item Ions of Z are positively charged.
\item Z has a greater electronegativity than V.
\item V has more outer electrons than Z.
\end{enumerate}
\sourcefoot{2023 FM Paper 12 Question 22}

% A3 typeset table — keep stem + table together
\Needspace{11em}
\qhead{A3.}{1}
{\small The elements in Period 3 and their compounds show trends across the period from sodium to chlorine.}

{\small Which row is correct?}

\vspace{0.12em}
\centering
{\small
\begin{tabular}{ccc}
\toprule
 & electronegativity of the elements & acid/base behaviour of the oxides of the elements \\
\midrule
\textbf{A} & decreases & basic $\rightarrow$ amphoteric $\rightarrow$ acidic \\
\textbf{B} & decreases & acidic $\rightarrow$ amphoteric $\rightarrow$ basic \\
\textbf{C} & increases & basic $\rightarrow$ amphoteric $\rightarrow$ acidic \\
\textbf{D} & increases & acidic $\rightarrow$ amphoteric $\rightarrow$ basic \\
\bottomrule
\end{tabular}}
\par\raggedright
\sourcefoot{2022 FM Paper 12 Question 18}

% A4 typeset table
\qhead{A4.}{1}
{\small Which row shows the expected properties of the element astatine when compared to the properties of the element iodine?}

\vspace{0.12em}
\centering
{\small
\begin{tabular}{ccc}
\toprule
 & electronegativity of astatine compared to iodine & volatility of astatine compared to iodine \\
\midrule
\textbf{A} & less electronegative & higher \\
\textbf{B} & more electronegative & higher \\
\textbf{C} & less electronegative & lower \\
\textbf{D} & more electronegative & lower \\
\bottomrule
\end{tabular}}
\par\raggedright
\sourcefoot{2023 MJ Paper 13 Question 17}

% A5 graphs
\clipq{A5.}{1}{img-005.png}
\sourcefoot{2025 MJ Paper 11 Question 20}

% A6 graphs NEW
\clipq{A6.}{1}{img-006.png}
\sourcefoot{2025 ON Paper 14 Question 17}

% A7 compact formulae
\qhead{A7.}{1}
{\small Which compound has the smallest difference in electronegativity between its two elements?}
\optline{\textbf{A} \ce{KF}\qquad \textbf{B} \ce{KBr}\qquad \textbf{C} \ce{LiF}\qquad \textbf{D} \ce{LiBr}}
\sourcefoot{2022 MJ Paper 13 Question 5}

% A8 typeset
\qhead{A8.}{1}
{\small L and M are elements in Period 3 of the Periodic Table. Neither element is argon.}

{\small Information about the Pauling electronegativity values of L and M is given.}

\vspace{0.12em}
\centering
{\small
\begin{tabular}{cc}
\toprule
element & Pauling electronegativity value \\
\midrule
L & the highest of the seven elements \ce{Na} to \ce{Cl} \\
M & the lowest of the seven elements \ce{Na} to \ce{Cl} \\
\bottomrule
\end{tabular}}
\par\raggedright

\vspace{0.1em}
{\small Three statements about elements L and M are given.}
\begin{enumerate}[label=\arabic*., leftmargin=1.6em, itemsep=0.06em]
\item Element L contains covalent bonds.
\item Element L has a higher atomic number than element M.
\item A compound of L and M contains ionic bonds.
\end{enumerate}
{\small Which statements are correct?}
\optline{\textbf{A} 1, 2 and 3\qquad \textbf{B} 1 and 2 only\qquad \textbf{C} 1 and 3 only\qquad \textbf{D} 2 and 3 only}
\sourcefoot{2023 MJ Paper 13 Question 4}

% A9 typeset — keep stem + both tables on one page
\Needspace{22em}
\qhead{A9.}{1}
{\small Elements X, Y and Z are all in the first two periods of the Periodic Table.}

{\small Their Pauling electronegativity values, $E_\mathrm{N}$, are shown.}

\vspace{0.1em}
\centering
{\small
\begin{tabular}{cc}
\toprule
element & $E_\mathrm{N}$ \\
\midrule
X & 1.0 \\
Y & 2.1 \\
Z & 4.0 \\
\bottomrule
\end{tabular}}
\par\raggedright

\vspace{0.08em}
{\small Substances exist with formulae \ce{XZ}, \ce{YZ} and \ce{Z2}.}

{\small Which row puts these substances in order of increasing melting point?}

\vspace{0.12em}
\centering
{\small
\begin{tabular}{cccc}
\toprule
 & lowest melting point & $\longrightarrow$ & highest melting point \\
\midrule
\textbf{A} & \ce{XZ} & \ce{YZ} & \ce{Z2} \\
\textbf{B} & \ce{XZ} & \ce{Z2} & \ce{YZ} \\
\textbf{C} & \ce{Z2} & \ce{YZ} & \ce{XZ} \\
\textbf{D} & \ce{Z2} & \ce{XZ} & \ce{YZ} \\
\bottomrule
\end{tabular}}
\par\raggedright
\sourcefoot{2022 FM Paper 12 Question 7}

% A10 NEW compact
\qhead{A10.}{1}
{\small Which type of bonding is \textbf{never} found in elements?}
\vspace{0.08em}
{\small
\begin{tabular}{@{}p{0.46\linewidth}p{0.46\linewidth}@{}}
\textbf{A} covalent & \textbf{B} ionic \\[0.15em]
\textbf{C} metallic & \textbf{D} van der Waals' forces
\end{tabular}}
\par
\sourcefoot{2021 MJ Paper 13 Question 6}

% A11 NEW
\qhead{A11.}{1}
{\small Sodium, magnesium, aluminium, silicon and phosphorus are all elements in Period 3 of the Periodic Table.}

{\small Three statements about the oxides and chlorides of these elements are given.}
\begin{enumerate}[label=\arabic*., leftmargin=1.6em, itemsep=0.06em]
\item The ionically bonded oxides all react with dilute hydrochloric acid.
\item All metal chlorides produce neutral solutions when added to water.
\item The two most electronegative elements both form covalently bonded chlorides.
\end{enumerate}
{\small Which statements are correct?}
\optline{\textbf{A} 1, 2 and 3\qquad \textbf{B} 1 and 2 only\qquad \textbf{C} 1 and 3 only\qquad \textbf{D} 2 and 3 only}
\sourcefoot{2023 MJ Paper 12 Question 19}

\vspace{0.35em}
{\large\bfseries Section B --- Structured questions (5 marks)}\par\vspace{0.1em}
{\normalsize\bfseries B1 · 2023 February/March Paper 22 Question 1(a) (3 marks)}\par

{\small The Pauling electronegativity values of elements can be used to predict the chemical properties of compounds.}

{\small Use the information in Table 1.1 to answer the following questions.}

\vspace{0.12em}
\centering
{\small
\begin{tabular}{lccccc}
\toprule
element & \ce{H} & \ce{Li} & \ce{C} & \ce{O} & \ce{S} \\
\midrule
Pauling electronegativity value & 2.1 & 1.0 & 2.5 & 3.5 & 2.6 \\
\bottomrule
\end{tabular}}
\par\vspace{0.06em}
{\small Table 1.1}\par
\raggedright

\qhead{B1(a)(i)}{1}
{\small Define electronegativity.}
\answerlines{2}

\qhead{B1(a)(ii)}{2}
{\small \ce{O} and \ce{S} are in Group 16.}

{\small Explain the difference in the Pauling electronegativity values of \ce{O} and \ce{S}.}
\answerlines{2}

\vspace{0.25em}
{\normalsize\bfseries B2 · 2023 May/June Paper 22 Question 1(b)(ii) (2 marks)}\par

\qhead{B2(b)(ii)}{2}
{\small Explain why electronegativity increases across a period.}
\answerlines{2}

\end{document}
